\documentclass[12pt]{article}
\usepackage{amsmath,amssymb}
\newcommand{\E}{\mathbb E}
\newcommand{\Law}{\operatorname{Law}}
\newcommand{\Ent}{\operatorname{Ent}}
\begin{document}
\section*{Research notes}

Problem: uniqueness of invariant probability for Dirichlet stochastic heat
equation on $[0,1]$ with additive space-time white noise and distributional
drift $\delta_0(u)$, under the ABLM regularized/mild interpretation and
regularisation by $G_{1/n}$.

Key interpretation:
\[
\delta_0(r)=-\frac12(-2{\bf1}_{r>0})',
\]
so the equation is a singular gradient SPDE with bounded potential
$f(r)=-2H_+(r)$.  For smooth approximation $B_n'=G_{1/n}$ with
$B_n(r)=\int_{-\infty}^rG_{1/n}(s)\,ds$, the invariant measure of the
regularized gradient equation is
\[
\pi_n(dv)=Z_n^{-1}\exp\left(2\int_0^1 B_n(v(x))\,dx\right)\gamma(dv),
\]
where $\gamma$ is Brownian bridge law (the invariant measure of the
Dirichlet OU/string process).  Since $0\le B_n\le1$ and $B_n\to H_+$
away from zero, while a Brownian bridge spends zero Lebesgue time at
zero,
\[
\pi_n\to
\pi(dv)=Z^{-1}\exp\left(2\int_0^1{\bf1}_{v(x)>0}\,dx\right)\gamma(dv)
\]
in total variation.  Densities $d\pi_n/d\gamma$ and $d\pi/d\gamma$ are
bounded above and below by deterministic positive constants.

Round-1 robust result obtained: assuming pointwise weak convergence of
regularized semigroups on bounded continuous observables,
\[
P_t^n\Phi(x)\to P_t\Phi(x),\qquad x\in C_0,\ t>0,
\]
one can pass reversibility/invariance and a uniform spectral gap from
$P^n$ to the actual ABLM semigroup $P$ without invoking
Bounebache--Zambotti:
\[
\int gP_tf\,d\pi=\int fP_tg\,d\pi,\qquad
\operatorname{Var}_\pi(P_tf)\le e^{-2ct}\operatorname{Var}_\pi(f).
\]
Key estimates:
\[
\operatorname{Var}_{\pi_n}(F)\le C\mathcal E_n(F,F),\qquad
\mathcal E_n(F,F)=\frac12\int\|DF\|_{L^2}^2\,d\pi_n,
\]
with $C$ uniform in $n$, by Gaussian Poincare for $\gamma$ plus bounded
density perturbation.  Then pass the $L^2$ contraction/spectral gap to
the limit using total variation convergence $\pi_n\to\pi$ and bounded
pointwise convergence $P_t^nf\to P_tf$.

Abstract uniqueness mechanism:
If $\eta$ is invariant and the actual resolvent satisfies
\[
R_\lambda^P(x,\cdot)\ll\pi\quad\hbox{for every }x,
\]
then $\eta=\lambda\eta R_\lambda^P\ll\pi$.  With $h=d\eta/d\pi$,
symmetry gives $P_th=h$ in $L^1(\pi)$.  For $h_a=h\wedge a$, concavity
gives $h_a\ge P_th_a$ and equal integrals force equality.  Since
$h_a\in L^2$, the spectral gap forces $h_a$ constant.  Varying $a$
implies $h$ constant, hence $\eta=\pi$.

Earlier route via Bounebache--Zambotti and its obstacle:
One possible route was to seek pointwise resolvent absolute continuity for the actual ABLM
semigroup, not only for the $L^2(\pi)$ Dirichlet-form version and not
only quasi-everywhere.  BZ Proposition 2.5 gives an $L^2$ resolvent kernel
$\rho_\lambda(x,dy)\ll\nu(dy)$ for all $x$ (setting it to zero on an
exceptional set) and transition absolute continuity for all
non-exceptional initial conditions.  BZ Proposition 3.1 constructs a
weak local-time solution for quasi-every initial condition.  ABLM
uniqueness plus the 2025 weak/mild equivalence paper plausibly identify
the BZ process with the ABLM mild process for quasi-every initial point,
but this still does not prove that arbitrary invariant measures cannot
charge the exceptional set, nor that $R_\lambda^P(x,\cdot)\ll\pi$ for
$x$ in that set.

Possible ways to make the Bounebache--Zambotti route work:
1. Uniform heat-kernel/capacity estimates for $P_t^n$: prove
   $P_t^n(x,A)\le C_{t,x}\,\mathrm{Cap}_n(A)^\alpha$ or
   $\|dP_t^n(x,\cdot)/d\pi_n\|_{L^2(\pi_n)}\le C_{t,x}$ with constants
   stable as $n\to\infty$ (perhaps only for $t$ large).  This would pass
   AC to the limit and imply uniqueness via convergence to equilibrium.
2. Strict version of BZ process: prove that the ABLM process is the
   strict all-starting-points Hunt process associated with the BZ
   Dirichlet form, and that BZ's resolvent kernel represents its
   resolvent at every deterministic $x$.
3. Entrance from exceptional sets: let $N$ be the BZ exceptional set.
   Show $P_s^P(x,N)=0$ for every $x$ and $s>0$.  Then q.e. resolvent AC
   plus the Markov property gives all-$x$ resolvent AC:
   $R_\lambda(x,A)\le\int_0^s e^{-\lambda t}P_t(x,A)dt$ for
   $\pi(A)=0$, then let $s\downarrow0$.
4. Prove uniqueness of invariant measures without pointwise AC by showing
   every invariant stationary solution has enough finite-energy/Sobolev
   regularity not to charge zero-capacity sets.  Stationarity and the
   nondegenerate additive noise imply one-dimensional projections have no
   atoms by the occupation formula, but this is far weaker than avoiding
   all Gaussian-capacity-zero sets.

Useful literature facts:
- BZ use $\nu(dx)\propto \exp(-\int f(x_r)dr)\mu(dx)$.  For
  $f=-2H_+$ this is exactly $\pi$.
- BZ equation has singular term
  $-\frac12\int f(da)\partial_\theta\ell^a_{t,\theta}$; for
  $f=-2H_+$ this is $\partial_\theta\ell^0$, formally
  $\delta_0(u(t,\theta))$.
- ABLM 2024 proves strong existence/uniqueness for $b=\kappa\delta_0$
  and convergence of smooth approximations; it explicitly did not treat
  Dirichlet boundary conditions, but the problem assumes the extension.
- ABLM 2025 ``Analytically weak and mild solutions...'' proves weak
  regularized and mild regularized solution notions coincide in the
  known existence range; may be useful for local-time-to-mild
  identification, though it still does not by itself remove the
  exceptional-set issue.



\section*{Additional reduction lemmas and strategy notes}

Strong-Feller/resolvent-Feller reduction.  Since $\pi$ is invariant and
has full support, if $A$ is Borel and $\pi(A)=0$, then
\[
  \int P_t{\bf1}_A\,d\pi=\pi(A)=0,\qquad
  \int R_\lambda{\bf1}_A\,d\pi=\lambda^{-1}\pi(A)=0.
\]
Thus either of the following immediately upgrades $\pi$-a.e. null-set
avoidance to all starting points:
\[
  P_t:B_b(E)\to C_b(E)\quad \hbox{for every }t>0,
\]
or
\[
  R_\lambda:B_b(E)\to C_b(E)\quad \hbox{for some }\lambda>0.
\]
Indeed a nonnegative continuous function with zero integral against a
full-support probability is identically zero.  This proves
$P_t(x,\cdot)\ll\pi$ for all $x,t>0$ in the strong-Feller case, and
$R_\lambda(x,\cdot)\ll\pi$ in the resolvent-Feller case.

Entrance reduction.  Let $N$ be a BZ exceptional set such that the ABLM
and BZ resolvents are identified for $y\notin N$.  If
\[
  P_s^P(x,N)=0,\qquad x\in E,\ s>0,
\]
then all-starting-point resolvent absolute continuity follows.  For
$\pi(A)=0$,
\[
 R_\lambda(x,A)
 =\int_0^\varepsilon e^{-\lambda t}P_t(x,A)\,dt
  +e^{-\lambda\varepsilon}P_\varepsilon R_\lambda(\cdot,A)(x).
\]
The first term is at most $\varepsilon$ and the second term is at most
$\lambda^{-1}P_\varepsilon(x,N)$, hence zero; let
$\varepsilon\downarrow0$.

ABLM method note.  The ABLM 2024 paper explicitly says that its proof
uses neither Girsanov nor Zvonkin transformations; therefore one should
not rely on an unstated global Zvonkin homeomorphism to transfer strong
Feller or transition equivalence from an OU process.  Any strong-Feller
argument would have to be proved separately, probably using Malliavin,
capacity, or stochastic-sewing/Krylov estimates.

OU comparison.  For the linear Dirichlet stochastic heat equation,
\[
  Z_t=S_t x+\int_0^t S_{t-s}\,dW_s,
\]
the law at time $t>0$ is equivalent to $\gamma$.  In the sine basis with
Dirichlet eigenvalues $\lambda_k=(k\pi)^2$, the covariance eigenvalues are
$q_k(t)=(1-e^{-\lambda_k t})/\lambda_k$ up to the convention-dependent
factor, while $\gamma$ has eigenvalues proportional to $\lambda_k^{-1}$.
The covariance ratio differs from $1$ by a summable sequence and the mean
$S_t x$ is in the Cameron--Martin space $H^1_0$, so Feldman--Hajek gives
equivalence.  Thus it is enough, in principle, to control the singular
drift contribution relative to this Gaussian law.

Uniform-density route.  A sufficient estimate for the regularized
semigroups is: for every $t>0$ and $x\in E$, there exist $p>1$ and
$C_{t,x}<\infty$ such that
\[
  \sup_n \left\|
    \frac{dP_t^n(x,\cdot)}{d\pi_n}
  \right\|_{L^p(\pi_n)}\le C_{t,x}.
\]
Together with $\pi_n\to\pi$ in total variation and weak convergence
$P_t^n(x,\cdot)\Rightarrow P_t(x,\cdot)$, this implies
$P_t(x,\cdot)\ll\pi$.  A capacity version
\[
  P_t^n(x,O)\le C_{t,x}\operatorname{Cap}_n(O)^\alpha
\]
uniformly in $n$ would also imply entrance from BZ exceptional sets,
because the forms have densities bounded above and below relative to
Brownian bridge and hence comparable capacities.

Naive Girsanov issue.  For fixed $n$, the regularized drift
$G_{1/n}(u)$ is bounded and $H$-valued, so pathwise Girsanov from the
linear equation gives absolute continuity.  The usual Novikov/entropy
term contains
\[
  \int_0^t\|G_{1/n}(u_s)\|_{L^2(0,1)}^2\,ds,
\]
which may blow as $n\to\infty$ because $G_{1/n}^2$ approximates a
multiple of $\sqrt n\,\delta_0$.  Any successful limiting argument must
use cancellation or estimates depending on the bounded potential
$0\le B_n\le1$, not on $\|G_{1/n}\|_\infty$ or $\|G_{1/n}\|_2$.

Dirichlet-form comparison idea.  The regularized and limiting invariant
densities are uniformly bounded above and below relative to Brownian
bridge.  Hence the forms
\[
  \mathcal E_n(F,F)=\frac12\int\|DF\|_{L^2}^2\,d\pi_n,\qquad
  \mathcal E_0(F,F)=\frac12\int\|DF\|_{L^2}^2\,d\gamma
\]
have comparable $1$-energies and capacities.  This gives spectral-gap
and compact-resolvent information uniformly, but by itself only produces
kernels and absolute continuity in the $L^2(\pi)$ or quasi-everywhere
sense.  The missing upgrade is a strict, all-starting-points statement.



\section*{Additional reductions for the exceptional-set gap}

Eventual equicontinuity/e-property route.  Suppose $P_t$ is already known
to be reversible with invariant $\pi$, and $P_t f\to\pi(f)$ in
$L^2(\pi)$ for every bounded Lipschitz $f$.  If for every $x\in E$ and
$\varepsilon>0$ there are a neighbourhood $U$ of $x$ and $t_0$ such that
\[
  |P_t f(y)-P_t f(x)|\le \varepsilon,\qquad y\in U,\quad t\ge t_0,
\]
then $P_t f(x)\to\pi(f)$ for every $x$.  The proof avoids compactness:
since $\pi$ has full support, $\pi(U)>0$, and
\[
 |P_t f(x)-\pi(f)|
 \le \varepsilon+
 \left(\pi(U)^{-1}\int_U |P_t f(y)-\pi(f)|^2\,\pi(dy)\right)^{1/2}.
\]
The second term goes to $0$.  Then any invariant law $\eta$ satisfies
$\eta(f)=\int P_t f\,d\eta\to\pi(f)$; bounded Lipschitz functions
separate probability measures on $E$.  This route would prove uniqueness
without all-starting-point absolute continuity.  The analytic missing
input is the eventual equicontinuity estimate.  Ordinary finite-time
Feller continuity is not enough, because the modulus may deteriorate as
$t\to\infty$.

Possible use of ABLM stability estimates.  In the ABLM uniqueness proof
for two solutions driven by the same noise and with the same initial
condition, the local estimate has the form
\[
\|z_t-P_{t-s}z_s\|\lesssim
D_{s,t}(t-s)^{1/2+\delta}+D_{s,t}|\log D_{s,t}|(t-s),
\]
where $D_{s,t}=\sup_{r\in[s,t]}\|z_r\|$.  For different initial data the
same local identity should still hold for the nonlinear drift part,
while $P_tz_0$ remains as an inhomogeneous term.  This suggests finite
time continuous dependence (hence weak Feller), perhaps by an Osgood
argument.  It does not immediately imply the eventual equicontinuity
needed above; one would need either a uniform-in-time modulus or an
asymptotic-strong-Feller estimate in which the heat semigroup
dissipation beats the logarithmic loss.

Capacity-smooth invariant measures route.  Let $N$ be the BZ exceptional
set on whose complement the resolvent kernel represents the actual
resolvent and is absolutely continuous with respect to $\pi$.  It is
enough to show that every invariant probability $\eta$ is smooth for the
BZ Dirichlet form, in particular $\eta(N)=0$.  A sufficient finite-energy
inequality is
\[
  \left|\int \varphi\,d\eta\right|
  \le C_\eta \mathcal E_1(\varphi,\varphi)^{1/2}
\]
for bounded nonnegative quasi-continuous $\varphi\in D(\mathcal E)$.
Indeed, applying this to equilibrium potentials gives
\[
  \eta(O)\le C_\eta\,\operatorname{Cap}(O)^{1/2}
\]
for open $O$, hence $\eta$ charges no zero-capacity set.  Once
$\eta(N)=0$, for any $\pi$-null set $A$,
\[
  \eta(A)=\lambda\int_{N^c}R_\lambda(y,A)\,\eta(dy)=0,
\]
so $\eta\ll\pi$ and the $L^2(\pi)$ spectral gap forces $\eta=\pi$.
A possible way to prove the finite-energy inequality would be a
stationary Krylov/capacity estimate for the regularized processes with
constants depending on the bounded primitives $B_n$ rather than on
$\|G_{1/n}\|_\infty$.

Local-time form of the drift.  Bounebache--Zambotti write the SPDE with
a local-time-in-space term
\[
 -\frac12\int_\mathbb R f(da)\,\partial_\theta \ell^a_{t,\theta}.
\]
For $f=-2H_+$, $df=-2\delta_0$, so this term is
$\partial_\theta\ell^0_{t,\theta}$, which is formally
$\delta_0(u(t,\theta))$ by the occupation density formula in the spatial
variable.  Thus their Dirichlet form matches the Gibbs measure
$\pi\propto\exp(2\int_0^1{\bf1}_{\{v>0\}})\,d\gamma$, but their Markov
solution is constructed only outside a properly exceptional set.

Finite-dimensional projection/entropy route.  For an invariant law
$\eta$, finite-dimensional projections onto the first $m$ Dirichlet
modes should have Lebesgue densities because the projected dynamics has
a nondegenerate Brownian martingale part.  This by itself is insufficient
for $\eta\ll\pi$ in infinite dimensions.  A useful target would be
uniform relative entropy
\[
  \sup_m H((\Pi_m)_\#\eta\mid(\Pi_m)_\#\pi)<\infty,
\]
which would imply $\eta\ll\pi$ by projective convergence of entropy.
The difficulty is to derive such a bound from the stationary martingale
problem with the distributional drift; the drift contribution is a
surface/local-time measure rather than an $L^2$ vector field.

Roughness decomposition idea.  At fixed $t>0$, the linear stochastic
convolution has a Gaussian law equivalent to Brownian bridge.  If the
drift convolution were an independent Cameron--Martin shift, absolute
continuity would follow immediately.  In reality it is adapted and
depends on the same noise.  A possible approach is to split the driving
noise into independent high-frequency or late-time pieces and show that
the singular drift contribution cannot cancel the nondegenerate Gaussian
roughness.  This would amount to a Malliavin or conditional-density
argument for the singular/local-time SPDE.



\section*{Round 4 refined reductions and obstacles}

The finite-energy/capacity shortcut must be formulated carefully.  An
inequality involving quasi-continuous versions
\[
  \left|\int \tilde\varphi\,d\eta\right|
  \le C_\eta \mathcal E_1(\varphi,\varphi)^{1/2}
\]
is not well-defined before knowing that $\eta$ does not charge exceptional
sets, because $\tilde\varphi$ is only defined up to capacity-zero sets.
A safer non-circular criterion is a capacity--Krylov estimate: for every
compact $K\subset E$, $0<\varepsilon<T$, and open $O\subset E$,
\[
 \sup_{x\in K}\mathbf E_x\int_\varepsilon^T{\bf1}_O(u_s)\,ds
 \le C_{K,\varepsilon,T}\operatorname{Cap}(O)^\alpha .
\]
Then any invariant $\eta$ avoids zero-capacity sets: choose compact $K$
with $\eta(K^c)<\delta$ and use stationarity to bound
\[
 (T-\varepsilon)\eta(O)
 \le C_{K,\varepsilon,T}\operatorname{Cap}(O)^\alpha
      +(T-\varepsilon)\delta .
\]
For $\operatorname{Cap}(N)=0$, outer regularity gives open
$O\supset N$ with the first term arbitrarily small for this fixed $K$,
then $\delta\downarrow0$.  This avoids the circular equilibrium-potential
argument.

Absolute continuity of BZ resolvent q.e. plus finite-time weak Feller is
not enough.  If $x_m\to x$ and $P_t(x_m,\cdot)\ll\pi$, weak convergence
of $P_t(x_m,\cdot)$ to $P_t(x,\cdot)$ preserves zero mass only for open
$\pi$-null sets, which are essentially useless since $\pi$ has full
support.  Weak limits of absolutely continuous measures may be singular.

For fixed-time absolute continuity, the mild decomposition
\[
  u_t=S_tu_0+Z_t+\int_0^t S_{t-s}\,\delta_0(u_s)\,ds
\]
suggests that the drift convolution is smoother than the Brownian-bridge
rough part $Z_t$.  However a random Cameron--Martin shift depending on
the Gaussian variable need not preserve absolute continuity: in an
abstract Wiener space, the map that removes one Gaussian coordinate is
of the form $w\mapsto w+F(w)$ with $F(w)$ Cameron--Martin-valued, yet its
image is supported on a hyperplane.  Thus one needs adaptedness,
invertibility, Malliavin nondegeneracy, or a time-splitting argument; mere
pathwise Cameron--Martin regularity of the drift convolution is
insufficient.

Stationary Fokker--Planck route.  Classical Bogachev--Roeckner type
regularity theorems for invariant measures of infinite-dimensional
diffusions can imply absolute continuity with respect to a Gaussian
reference measure when the perturbing drift is Cameron--Martin/noise-Hilbert
valued and square integrable with respect to the invariant measure.  The
present drift $\delta_0(u(\cdot))$ is a spatial local-time measure, not an
$L^2(0,1)$ vector field; the natural Gibbs density is bounded but need
not have finite Fisher information for the $L^2$-gradient Dirichlet form.
Thus those theorems do not apply directly without a new argument adapted
to bounded-variation potentials.

A potentially useful invariant-measure-only strategy is to start the
regularized equation from an arbitrary invariant $\eta$ of the limiting
semigroup.  Then $\eta P_t^n\Rightarrow\eta$ for each fixed $t$ by (A).
If one had a uniform-in-$n$ entropy, $L^p$-density, or capacity estimate
for $\eta P_t^n$ relative to $\pi_n$ on compact subsets of initial data,
then $\eta\ll\pi$ would follow.  This is equivalent in strength to the
capacity--Krylov estimate above.  Naive Girsanov or Bismut constants for
$G_{1/n}$ blow up; any successful estimate must use the bounded primitive
$0\le B_n\le1$ rather than $\|G_{1/n}\|_\infty$.



\section*{All-starting fixed-time absolute continuity route}

A way to close the exceptional-set gap is to use the local-time form of
the ABLM drift and a time-shifted Girsanov argument.

For a deterministic start $x\in E$, write the ABLM solution as
\[
  u_t=S_tx+V_t+K_t,\qquad V_t=\int_0^tS_{t-s}\,dW_s .
\]
For non-negative measure drifts ABLM's random-control construction
identifies the distributional drift at time $s$ with a finite positive
spatial local-time measure
\[
  \mu_s(d\theta)=d_\theta L_s^0(\theta),
\]
where $L_s^a(\theta)$ is the local time at level $a$ of the spatial path
$\theta\mapsto u_s(\theta)$.  The mild drift is then
\[
  K_t=\int_0^tS_{t-s}\mu_s\,ds .
\]
At fixed $s>0$, the stochastic convolution $V_s$ has covariance
$(-\Delta_D)^{-1}(I-e^{s\Delta_D})$; in the sine basis it differs from a
Brownian bridge by a smooth Gaussian function.  ABLM estimates give that
$S_sx+K_s$ is spatially $C^\alpha$ for every $\alpha<1$, locally uniformly
for $s>0$.  Hence $u_s(\cdot)$ is Brownian-bridge-like plus a zero
quadratic variation perturbation and has a spatial local time satisfying
the occupation formula
\[
  \int_0^1 \varphi(\theta)g(u_s(\theta))\,d\theta
  =\int_{\mathbb R}g(a)\left(\int_0^1\varphi(\theta)\,d_\theta
     L_s^a(\theta)\right)\,da .
\]
Thus $G_\varepsilon(u_s(\theta))\,d\theta\to\mu_s(d\theta)$ weakly.

The key energy estimate needed for terminal absolute continuity is
\[
  \int_0^t\|S_{(t-s)/2}\mu_s\|_{L^2(0,1)}^2\,ds<\infty\quad\text{a.s.}
\]
It follows from the heat-kernel bound
\[
  \|S_r\nu\|_2^2=\iint p_{2r}(a,b)\,\nu(da)\nu(db)
     \le C r^{-1/2}\|\nu\|_{\mathrm{TV}}^2
\]
and the ABLM local-time/Krylov moment estimate, for $m\ge2$,
\[
  \mathbb E[\|\mu_s\|_{\mathrm{TV}}^m]\le C_{m,T}(1+s^{-m/4}),\qquad
  0<s\le T .
\]
The singularities $s^{-1/2}$ and $(t-s)^{-1/2}$ are both integrable in
the energy expectation.

Time-shifted Girsanov algebra.  Suppose a mild equation has a predictable
finite-measure drift $F_s$ satisfying
\[
  \int_0^t\|S_{(t-s)/2}F_s\|_2^2\,ds<\infty .
\]
Set
\[
  \widehat F_s=2{\bf1}_{[t/2,t]}(s)S_{t-s}F_{2s-t}.
\]
Then $\widehat F_s$ is $L^2$-valued and adapted because $2s-t\le s$ on
$[t/2,t]$, and
\[
  \int_0^t\|\widehat F_s\|_2^2ds
  =2\int_0^t\|S_{(t-r)/2}F_r\|_2^2dr<\infty .
\]
The auxiliary linear equation
\[
  \widehat u_r=S_rx+V_r+\int_0^rS_{r-s}\widehat F_s\,ds
\]
has path law absolutely continuous with respect to the OU law by the
ordinary Hilbert-space Girsanov theorem, after standard localization.
At the terminal time,
\[
  \int_0^tS_{t-s}\widehat F_s\,ds
  =\int_0^tS_{t-r}F_r\,dr,
\]
so $\widehat u_t$ equals the original terminal variable.  Therefore
$P_t(x,\cdot)$ is absolutely continuous with respect to the OU one-time
law for every deterministic $x$.  The OU law at time $t>0$ is equivalent
to Brownian bridge measure since its covariance eigenvalues are
$(1-e^{-k^2\pi^2t})/(k^2\pi^2)$ and the mean $S_tx$ lies in $H_0^1$.
Consequently $P_t(x,\cdot)\ll\gamma\sim\pi$ for all $x$, which implies
that any invariant measure is absolutely continuous with respect to
$\pi$ and hence equals $\pi$ by the spectral-gap argument.



\section*{Dyadic adapted-control proof of fixed-time absolute continuity}

The local-time route can be avoided.  The only ABLM input needed for
all-starting fixed-time absolute continuity is the $V(3/4)$ drift estimate.
For a deterministic start write
\[
  u_t=S_t x+V_t+K_t,\qquad K_0=0,
\]
and use the ABLM measure-drift regularity
\[
 \sup_{0<s<t\le T}\sup_{\theta\in[0,1]}
 \frac{\|K_t(\theta)-S_{t-s}K_s(\theta)\|_{L^m}}{|t-s|^{3/4}}<\infty .
\]
Fix a terminal time $t$ and set $r_k=t(1-2^{-k})$, $\Delta_k=r_{k+1}-r_k$.
Let
\[
  D_k=K_{r_{k+1}}-S_{\Delta_k}K_{r_k}.
\]
Then $D_k$ is $\mathcal F_{r_{k+1}}$-measurable and, for $k\ge1$,
\[
  \mathbb E\|D_k\|_{L^2}^2\lesssim \Delta_k^{3/2}.
\]
Define a control on the next dyadic interval:
\[
  F_s=\sum_{k\ge0} {\bf1}_{(r_{k+1},r_{k+2}]}(s)\,
       \Delta_{k+1}^{-1}S_{s-r_{k+1}}D_k .
\]
This is adapted/predictable because the increment $D_k$ is known by the
left endpoint $r_{k+1}$ of the interval on which it is used.  Its energy is
finite:
\[
  \mathbb E\int_0^t\|F_s\|_2^2\,ds
  \le \frac{\mathbb E\|D_0\|_2^2}{\Delta_1}
     +\sum_{k\ge1}\frac{\mathbb E\|D_k\|_2^2}{\Delta_{k+1}}
  \lesssim 1+\sum_{k\ge1}\Delta_k^{1/2}<\infty .
\]
At the terminal time the control reproduces exactly the singular drift:
\[
 \int_0^tS_{t-s}F_s\,ds
 =\sum_{k\ge0}S_{t-r_{k+1}}D_k
 =\lim_{N\to\infty}S_{t-r_{N+1}}K_{r_{N+1}}
 =K_t .
\]
The sum telescopes.  Therefore
\[
 \widehat u_s=S_sx+V_s+\int_0^sS_{s-r}F_r\,dr
\]
has $\widehat u_t=u_t$.  Since $F$ is adapted and has a.s. finite
$L^2([0,t];H)$ energy, localized Hilbert-space Girsanov gives
\[
  \mathcal L(u_t)\ll \mathcal L(S_tx+V_t).
\]
The OU one-time law is equivalent to Brownian bridge by Feldman--Hajek:
in the sine basis the covariance eigenvalues are
$(1-e^{-k^2\pi^2 t})/(k^2\pi^2)$ compared with $(k^2\pi^2)^{-1}$, and
the mean $S_tx$ is in $H_0^1$.  Hence $P_t(x,\cdot)\ll\gamma\sim\pi$ for
every deterministic $x$ and $t>0$.

This route is robust because the exponent $3/4>1/2$ is exactly what makes
the dyadic energy series $\sum\Delta_k^{2(3/4)-1}$ converge.  More
generally any adapted mild drift $K$ with parabolic increments of exponent
$\alpha>1/2$ admits such a terminal adapted-control representation.


\section*{Dirichlet boundary compensation for the ABLM \(V(3/4)\) estimate}

The dyadic adapted-control proof only needs the parabolic drift increment
estimate
\[
  \|K_t(x)-S_{t-s}K_s(x)\|_{L^m}\lesssim |t-s|^{3/4}.
\]
ABLM prove this for the line, periodic, and Neumann kernels.  For the
Dirichlet kernel the pointwise local nondeterminism of the stochastic
convolution degenerates near \(0,1\), but the heat kernel also vanishes
linearly at the boundary, and the two effects compensate.

Let
\[
  \rho_\tau(y)=\int_0^\tau p^D_{2a}(y,y)\,da,\qquad d(y)=y\wedge(1-y).
\]
For the half-line kernel \(p^{+}_a(y,y)=g_a(0)-g_a(2y)\),
\[
  \int_0^\tau p^{+}_{2a}(y,y)\,da\asymp \sqrt{\tau}\wedge y .
\]
By reflection and localization the interval kernel satisfies, uniformly on
bounded time intervals,
\[
  \rho_\tau(y)\asymp \sqrt{\tau}\wedge d(y).
\]
Thus the conditional density of
\(V_{s+\tau}(y)-S_\tau V_s(y)\) at any point is bounded by
\(C(\sqrt{\tau}\wedge d(y))^{-1/2}\), and the derivative bound for the
one-dimensional Gaussian smoothing of a unit finite measure is bounded by
\(C(\sqrt{\tau}\wedge d(y))^{-1}\).

The required weighted integrals are
\[
 \sup_x\int_0^1p^D_a(x,y)(\sqrt{\tau}\wedge d(y))^{-1/2}\,dy
 \lesssim a^{-1/4}+\tau^{-1/4},
\]
and
\[
 \sup_x\int_0^1p^D_a(x,y)(\sqrt{\tau}\wedge d(y))^{-1}\,dy
 \lesssim a^{-1/2}+\tau^{-1/2}.
\]
For the boundary part, use the half-line scaling
\(p_a^+(x,y)=a^{-1/2}p_1^+(x/\sqrt a,y/\sqrt a)\).  Since
\(p_1^+(\xi,z)=O(z)\) as \(z\downarrow0\),
\[
 \sup_x\int_0^1p^D_a(x,y)d(y)^{-1/2}\,dy\lesssim a^{-1/4},\qquad
 \sup_x\int_0^1p^D_a(x,y)d(y)^{-1}\,dy\lesssim a^{-1/2}.
\]
The interior part gives the \(\tau^{-1/4}\) and \(\tau^{-1/2}\) terms.

Consequently the ABLM estimate corresponding to their Lemma 5.1 remains
valid for a non-negative unit-mass mollifier \(b_n\):
\[
 \left\|\int_s^t\!\!\int_0^1p^D_{T-r}(x,y)
 b_n(V_r(y)+S_{r-s}\kappa(y))\,dy\,dr\right\|_{L^m}
 \lesssim (t-s)^{3/4},
\]
for every \({\cal F}_s\)-measurable shift \(\kappa\).  For integer moments,
expand the \(m\)-th power, order the times, and condition successively; each
conditioning contributes the density factor above with
\(\tau=r_i-r_{i-1}\), and the ordered simplex integral of
\(\prod[(r_i-r_{i-1})^{-1/4}+(T-r_i)^{-1/4}]\) scales like
\((t-s)^{3m/4}\).  Non-integer \(m\) follows by Jensen/interpolation.

The conditional sewing increment in ABLM Proposition 3.8 uses the derivative
version and gives
\[
  |\mathbb E_u\delta A^{T,n}_{s,u,t}(x)|
   \lesssim \lambda^T_{s,u}(x)(t-u)^{1/2},
\]
where
\(\lambda^T_{s,u}(x)=S_{T-u}K_u(x)-S_{T-s}K_s(x)\ge0\) is the random
control coming from positivity of the measure drift.  Therefore ABLM's
stochastic sewing argument with random controls carries over to Dirichlet
boundary conditions for finite non-negative measure drifts.  This supplies
the \(V(3/4)\) estimate needed by the terminal Girsanov argument without
appealing to Bounebache--Zambotti exceptional-set potential theory.



\section*{Hilbert-space replacement for the Dirichlet \(V(3/4)\) issue}

The pointwise Dirichlet \(V(3/4)\) drift estimate is vulnerable near the
boundary: a weighted bridge inequality with factor
\(\rho_\tau(y)^{-1}\simeq(\sqrt{\tau}\wedge d(y))^{-1}\) fails if one tries to
factor a positive function through \(S_{T-u}\) pointwise.  For the terminal
Girsanov argument, however, pointwise control is stronger than necessary.  It
is enough to have the Hilbert estimate
\[
 \|K_t-S_{t-s}K_s\|_{L^2(\Omega;L^2(0,1))}
       \lesssim |t-s|^{3/4}.
\]

For the regularized equation \(u^n=S_tx+V_t+K_t^n\), set
\(D^n_{s,t}=K_t^n-S_{t-s}K_s^n\) and define the frozen approximation
\[
 A^n_{s,t}=\int_s^t S_{t-r}b_n(S_{r-s}u_s^n+V_r-S_{r-s}V_s)\,dr .
\]
The basic Dirichlet variance is
\[
 \rho_\tau(y)=\int_0^\tau p_{2a}(y,y)\,da
   \asymp \sqrt{\tau}\wedge d(y),\qquad d(y)=y\wedge(1-y).
\]
Consequences:
\[
 \int \rho_\tau^{-1/2}\lesssim \tau^{-1/4},\qquad
 \sup_z\int p_a(z,y)\rho_\tau(y)^{-1/2}\,dy
       \lesssim a^{-1/4}+\tau^{-1/4}.
\]
Sequential Gaussian conditioning gives
\[
 \|A^n_{s,t}\|_{L^2(\Omega;H)}\lesssim |t-s|^{3/4}
\]
uniformly in \(n\).  In detail, after ordering times \(s<r<q<t\),
\[
\E_s b_n(Z_r(y))b_n(Z_q(z))
 \lesssim \rho_{r-s}(y)^{-1/2}\rho_{q-r}(z)^{-1/2},
\]
and the \(H\)-inner product of heat kernels produces
\(p_{2t-r-q}(y,z)\); the resulting two-time integral scales as
\(|t-s|^{3/2}\).

The sewing defect is controlled in \(H\), not pointwise.  For \(s<u<t\),
the two frozen arguments for times \(r>u\) differ by \(S_{r-u}D^n_{s,u}\).
Conditional on \(\mathcal F_u\), the future Gaussian variance at \(y\) is
\(\rho_{r-u}(y)\), so for the Gaussian mollifiers
\[
 \left|\E_u\{b_n(Z_r(y))-b_n(Z_r(y)+h(y))\}\right|
   \lesssim \rho_{r-u}(y)^{-1}h(y),\qquad h=S_{r-u}D^n_{s,u}.
\]
The correct Dirichlet boundary compensation is Hardy's inequality:
\[
 \|\rho_\tau^{-1}S_\tau f\|_{2}
 \lesssim \tau^{-1/2}\|f\|_2.
\]
Indeed \(\rho_\tau^{-1}\lesssim \tau^{-1/2}+d^{-1}\), while
\(\|S_\tau f/d\|_2\lesssim\|\partial_xS_\tau f\|_2
 \lesssim \tau^{-1/2}\|f\|_2\).  Hence
\[
 \|\E_u\delta A^n_{s,u,t}\|_H
   \lesssim (t-u)^{1/2}\|D^n_{s,u}\|_H.
\]
Applying the Hilbert-valued stochastic sewing lemma on intervals of length
\(h\), with
\(Y_n(h)=\sup_{t-s\le h}\|D^n_{s,t}\|_{L^2(H)}/|t-s|^{3/4}\), gives
\(Y_n(h)\le C+C h^{1/2}Y_n(h)\).  For \(h\) small this is uniform in \(n\);
larger intervals follow by semigroup additivity.  This supplies the energy
estimate needed for the adapted dyadic control \(F^n\) in the terminal
Girsanov proof, without using the false pointwise bridge inequality.  The
one-dimensional derivative estimate must contain an absolute value:
\[
 \left|\E_u\{b_n(Z_r(y))-b_n(Z_r(y)+h(y))\}\right|
   \lesssim \rho_{r-u}(y)^{-1}|h(y)| .
\]

\section*{Entropy lower-semicontinuity route for terminal absolute continuity}

A stronger convergence \(K^n\to K\) of the drift additive functionals is not
needed for invariant-measure uniqueness.  For each fixed \(n\), the uniform
Hilbert estimate gives a dyadic adapted terminal control on \([0,t]\):
\[
 r_k=t(1-2^{-k}),\quad \Delta_k=r_{k+1}-r_k,\quad
 D_k^n=K^n_{r_{k+1}}-S_{\Delta_k}K^n_{r_k},
\]
\[
 F_s^n=\sum_{k\ge0}\frac{{\bf1}_{(r_{k+1},r_{k+2}]}(s)}
             {\Delta_{k+1}}S_{s-r_{k+1}}D_k^n .
\]
Then \(F^n\) is predictable and
\[
 \E\int_0^t\|F_s^n\|_2^2\,ds
 \lesssim \sum_k\Delta_k^{3/2}/\Delta_{k+1}<\infty
\]
uniformly in \(n\).  The terminal effect telescopes:
\[
 \int_0^tS_{t-s}F_s^n\,ds
 =\sum_{k\ge0}S_{t-r_{k+1}}D_k^n
 =K_t^n .
\]
Therefore, with \(\nu_t^x=\Law(S_tx+V_t)\), Hilbert-space Girsanov plus
data processing gives
\[
 \Ent(P_t^n(x,\cdot)\mid\nu_t^x)
 \le \frac12\E\int_0^t\|F_s^n\|_2^2\,ds\le C_t .
\]
Truncate \(F^n\) by its accumulated energy to justify Novikov; remove the
truncation by lower semicontinuity of relative entropy and convergence of the
controlled terminal values.

The stipulated semigroup convergence
\(P_t^n(x,\cdot)\Rightarrow P_t(x,\cdot)\) on \(E=C_0([0,1])\), together with
lower semicontinuity of relative entropy on Polish spaces, yields
\[
 \Ent(P_t(x,\cdot)\mid\nu_t^x)\le C_t,
\]
so \(P_t(x,\cdot)\ll\nu_t^x\).  Feldman--Hajek gives
\(\nu_t^x\sim\gamma\): in the Dirichlet sine basis the covariance eigenvalues
are \((1-e^{-k^2\pi^2t})/(k^2\pi^2)\), equivalent to the Brownian-bridge
covariance \((k^2\pi^2)^{-1}\), and \(S_tx\in H_0^1\).  Hence
\(P_t(x,\cdot)\ll\gamma\sim\pi\) for every deterministic initial condition,
which is the absolute-continuity input needed to prove uniqueness of the
invariant probability measure.



\section*{Consolidated proof details for the final argument}

The state space is $E=C_0([0,1])$ and $H=L^2(0,1)$.  The approximation
assumption is used as pointwise convergence of the regularized semigroups
$P_t^n\Phi(x)\to P_t\Phi(x)$ for all $\Phi\in C_b(E)$, $x\in E$, $t>0$.

For the regularized equations $b_n=G_{1/n}$, write
\[
u_t^n=S_tx+V_t+K_t^n,\qquad
K_t^n=\int_0^tS_{t-r}b_n(u_r^n)\,dr,\qquad
D^n_{s,t}=K_t^n-S_{t-s}K_s^n .
\]
The key estimate needed for endpoint absolute continuity is
\[
\|D^n_{s,t}\|_{L^2(\Omega;H)}\lesssim_T |t-s|^{3/4}
\]
uniformly in $n$ and $x$.

Useful heat-kernel facts.  With
\[
\rho_\tau(y)=\int_0^\tau p_{2a}(y,y)\,da,\qquad d(y)=y\wedge(1-y),
\]
the one-dimensional Dirichlet heat-kernel bound
$p_a(y,y)\asymp a^{-1/2}(1\wedge d(y)^2/a)$ implies
$\rho_\tau(y)\asymp \sqrt{\tau}\wedge d(y)$.  Consequently
\[
\int\rho_\tau^{-1/2}\lesssim \tau^{-1/4},\qquad
\sup_z\int p_a(z,y)\rho_\tau(y)^{-1/2}dy
  \lesssim a^{-1/4}+\tau^{-1/4},
\]
and
\[
\|\rho_\tau^{-1}S_\tau f\|_2\lesssim \tau^{-1/2}\|f\|_2.
\]
The last estimate follows from $\rho_\tau^{-1}\lesssim\tau^{-1/2}+d^{-1}$,
Hardy's inequality, and $\|\partial_xS_\tau f\|_2\lesssim\tau^{-1/2}\|f\|_2$.

Uniform Gaussian convolution estimates for $b_n=G_{1/n}$:
if $X=m+N(0,\sigma^2)$, then
$\sup_m\E b_n(X)\le (2\pi\sigma^2)^{-1/2}$ uniformly in $n$ because $b_n$ is
a probability density.  Also
\[
\sup_m |(b_n*\varphi_\sigma)'(m)|
   \le \sup_x|\varphi_\sigma'(x)|\lesssim\sigma^{-2},
\]
so
\[
|\E\{b_n(X+h)-b_n(X)\}|\lesssim \sigma^{-2}|h|.
\]

Frozen increments:
\[
A^n_{s,t}=\int_s^tS_{t-r}
 b_n(S_{r-s}u_s^n+V_r-S_{r-s}V_s)\,dr .
\]
For fixed $n$, dyadic Riemann sums of $A^n$ converge to $D^n$: compare on
$[a,b]$ with the actual integral and use
$\|D^n_{a,r}\|_2\le \|b_n\|_\infty(r-a)$ and $\operatorname{Lip}(b_n)<\infty$,
giving total error $O_n((t-s)|\Pi|)$.

For $Z_r=S_{r-s}u_s^n+V_r-S_{r-s}V_s$, conditional Gaussian estimates yield
\[
\E_s b_n(Z_r(y))b_n(Z_q(z))
\lesssim
\rho_{r-s}(y)^{-1/2}\rho_{q-r}(z)^{-1/2}.
\]
Combining with the heat-kernel convolution identity gives
$\|A^n_{s,t}\|_{L^2H}\lesssim |t-s|^{3/4}$.  The semigroup coboundary
$\delta A_{s,u,t}=A_{s,t}-S_{t-u}A_{s,u}-A_{u,t}$ satisfies
\[
\|\E_u\delta A^n_{s,u,t}\|_2
 \lesssim (t-u)^{1/2}\|D^n_{s,u}\|_2,
\]
because the frozen arguments after $u$ differ by $S_{r-u}D^n_{s,u}$ and the
last heat-kernel estimate controls
$\|\rho_{r-u}^{-1}S_{r-u}D^n_{s,u}\|_2$.

The semigroup stochastic sewing lemma in $L^2(H)$ can be proved directly from
dyadic refinements.  The difference between two consecutive dyadic sums is a
sum of $S_{t-b}\delta A_{a,u,b}$ terms.  Conditional means are controlled by a
geometric series when their exponent is $>1$; centered parts are orthogonal
Hilbert-valued martingale differences and are square-summable when the
centered exponent is $>1/2$.  Applying this with exponents $5/4$ and $3/4$
and bootstrapping
\[
Y_n(h)=\sup_{|t-s|\le h}\frac{\|D^n_{s,t}\|_{L^2H}}{|t-s|^{3/4}}
\]
gives $Y_n(h)\le C+Ch^{1/2}Y_n(h)$, hence the uniform $3/4$ estimate.

The Gibbs measures are
\[
\pi_n(dv)=Z_n^{-1}\exp\left(2\int_0^1B_n(v(r))\,dr\right)\gamma(dv),\quad
\pi(dv)=Z^{-1}e^{2\int_0^1{\bf1}_{v(r)>0}dr}\gamma(dv).
\]
Total variation convergence follows from bounded convergence and the fact
that Brownian bridge spends zero Lebesgue time at level $0$.  Regularized
reversibility can be justified via cylinder-potential Galerkin
approximations: replace $2\int B_n(v)$ by $2\int B_n(\Pi_m v)$, obtaining a
finite-dimensional gradient perturbation of the reversible OU process plus an
independent OU tail.  Let $m\to\infty$ using Lipschitz convergence of the mild
solutions.  Holley--Stroock gives a Poincare constant for $\pi_n$ uniform in
$n$ because the densities relative to $\gamma$ are bounded above and below
uniformly; semigroup convergence and $\pi_n\to\pi$ pass invariance, symmetry
and the spectral gap to $P_t$.

Endpoint absolute continuity: choose $r_k=t(1-2^{-k})$ and define the adapted
control
\[
F_s^n=\sum_{k\ge0}\Delta_{k+1}^{-1}{\bf1}_{(r_{k+1},r_{k+2}]}(s)
      S_{s-r_{k+1}}D^n_{r_k,r_{k+1}} .
\]
Then $\E\int_0^t\|F_s^n\|_2^2ds\lesssim
\sum\Delta_k^{3/2}/\Delta_{k+1}<\infty$ uniformly in $n$ and
\[
\int_0^tS_{t-s}F_s^n\,ds=\sum_{k\ge0}S_{t-r_{k+1}}D^n_{r_k,r_{k+1}}=K_t^n .
\]
After energy truncation, Hilbert-space Girsanov gives
\[
\Ent(\Law(S_tx+V_t+\int_0^tS_{t-s}F_s^{n,N}ds)\mid\Law(S_tx+V_t))
 \le \frac12\E\int_0^t\|F_s^n\|_2^2ds.
\]
The controlled variables are $E$-valued because parabolic smoothing maps
$L^2(0,t;H)$ into $H^\alpha_0\subset E$ for any $\alpha\in(1/2,1)$.  Let
$N\to\infty$ and then $n\to\infty$; lower semicontinuity of entropy on $E$
gives $P_t(x,\cdot)\ll\nu_t^x=\Law(S_tx+V_t)$.  Feldman--Hajek on $E$ gives
$\nu_t^x\sim\gamma$ because the covariance eigenvalue ratios are
$1-e^{-k^2\pi^2t}$ and $S_tx\in H^1_0$.

Finally, if $\eta$ is invariant, fixed-time absolute continuity implies
$\eta\ll\pi$, say $d\eta=h\,d\pi$.  Symmetry gives $P_th=h$ in $L^1(\pi)$.
For $h_a=h\wedge a$, Jensen gives $P_th_a\le h_a$, equal integrals force
equality, and the spectral gap implies $h_a$ is constant.  Letting
$a\uparrow\infty$ yields $h\equiv1$.


\section*{Additional technical details for a complete uniqueness proof}

For the weighted heat-kernel estimate, it is enough to use
\[
p_a(z,y)\le Ca^{-1/2}e^{-|z-y|^2/(Ca)}
\]
and $d(y)^{-1/2}\le y^{-1/2}+(1-y)^{-1/2}$.  The first endpoint term satisfies
\[
a^{-1/2}\int_0^1e^{-(y-z)^2/(Ca)}y^{-1/2}\,dy
 =a^{-1/4}\int_0^{1/\sqrt a}e^{-(r-z/\sqrt a)^2/C}r^{-1/2}\,dr
 \le Ca^{-1/4},
\]
because
$\sup_{\alpha\ge0}\int_0^\infty e^{-(r-\alpha)^2/C}r^{-1/2}dr<\infty$.
The second endpoint is identical.  Together with
$\rho_\tau\simeq \sqrt{\tau}\wedge d$ this yields
\[
\sup_z\int p_a(z,y)\rho_\tau(y)^{-1/2}dy
 \lesssim a^{-1/4}+\tau^{-1/4}.
\]
The estimate
\[
\|\rho_\tau^{-1}S_\tau f\|_2\lesssim \tau^{-1/2}\|f\|_2
\]
uses $\rho_\tau^{-1}\lesssim \tau^{-1/2}+d^{-1}$ and Hardy:
$\|g/d\|_2\lesssim \|g'\|_2$ for $g\in H^1_0$.

The semigroup sewing lemma can be recorded in the following exact dyadic
form.  For $J_m=\sum_{i=0}^{2^m-1}S_{t-t_{i+1}}A_{t_i,t_{i+1}}$ on the equal
dyadic partition, one has
\[
J_{m+1}-J_m=-\sum_i S_{t-t_{i+1}}\delta A_{t_i,u_i,t_{i+1}}.
\]
If
$\|\E_{t_i}\delta A_{t_i,u_i,t_{i+1}}\|_{L^2H}\le \Gamma_1\ell_m^{1+\epsilon_1}$
and
$\|\delta A_{t_i,u_i,t_{i+1}}-\E_{t_i}\delta A_{t_i,u_i,t_{i+1}}\|_{L^2H}
\le \Gamma_2\ell_m^{1/2+\epsilon_2}$, then
\[
\|\E\text{-mean part}\|_{L^2H}\le
\Gamma_1 |t-s|^{1+\epsilon_1}2^{-m\epsilon_1},
\]
and the centered summands are pairwise orthogonal, hence
\[
\|\text{centered part}\|_{L^2H}\le
\Gamma_2 |t-s|^{1/2+\epsilon_2}2^{-m\epsilon_2}.
\]
Summing over $m$ gives the sewing bound.

For regularized reversibility, the useful Galerkin approximation is not the
usual spectral truncation of the whole equation but the cylinder-potential
gradient equation
\[
u_t^{n,m}=S_tx+V_t+\int_0^tS_{t-r}\Pi_m b_n(\Pi_m u_r^{n,m})\,dr.
\]
It is reversible for
\[
\pi_{n,m}(dv)=Z_{n,m}^{-1}\exp\{2\int_0^1B_n(\Pi_m v(r))\,dr\}\gamma(dv),
\]
because the drift equals the $H$-gradient of the cylinder potential.  For
fixed $n$, convergence to the full regularized equation follows from
\[
\Pi_m b_n(\Pi_m z)-b_n(y)
=\Pi_m[b_n(\Pi_m z)-b_n(\Pi_m y)]
 +\Pi_m[b_n(\Pi_m y)-b_n(y)]+(\Pi_m-I)b_n(y),
\]
Lipschitz continuity of the Nemytskii map $b_n:H\to H$, dominated convergence
for the last two terms, and Gronwall.  At a fixed positive terminal time,
$E$-convergence follows by splitting the drift convolution at $t-\varepsilon$:
on $[0,t-\varepsilon]$ the heat semigroup maps $H$ errors into
$H_0^\alpha\subset C_0$ for $\alpha>1/2$, while the contribution of
$[t-\varepsilon,t]$ is uniformly small since $b_n$ is bounded.

The final use of symmetry with an invariant density $h=d\eta/d\pi$ only needs
$L^1$ duality.  Since $\pi$ is invariant, $P_t$ is a contraction on
$L^1(\pi)$ and $L^\infty(\pi)$.  For $g\in L^1$ and bounded $f$, truncate
$g_M=(-M)\vee(g\wedge M)$ and use $L^2$ symmetry:
$\int f P_tg_M\,d\pi=\int P_tf\,g_M\,d\pi$.  Letting $M\to\infty$ gives
$\int fP_tg\,d\pi=\int P_tf\,g\,d\pi$.  This justifies $P_th=h$ in
$L^1(\pi)$ for invariant $\eta$.

\end{document}



